This section states in one place what the twenty-two decisions above amount
to: exactly which inputs Phase~1 accepts, and what it reports for each input
it does not. It is the specification the implementation is written against,
and the acceptance condition is also the ingest-time validator.

Acceptance is a property of a \emph{pair}: an input and a configuration.
Nothing here is a property of the input alone.

\subsection*{A1. Inputs and configuration}

The input is a single JSON value, and must be either

\begin{itemize}
  \item an \textbf{object}, read as a corpus of one document; or
  \item an \textbf{array}, read as a corpus of its elements, every one of
        which must be an object.
\end{itemize}

\noindent
Any other top-level value is rejected, as is an array containing a non-object
element. A document's root is therefore always an object (S0). An empty array
is accepted and yields an empty schema.

The configuration supplies two path-directed markings, both optional:

\begin{lstlisting}[style=json]
maps      = [".users", ".index.by_id"]
recursive = [ { path = ".children[]", folds_into = "." } ]
\end{lstlisting}

\noindent
\texttt{maps} names object-valued paths whose members are data rather than
schema (S9). \texttt{recursive} identifies a path with an ancestor path it
folds into (S19).

\subsection*{A2. The type lattice}

Inference assigns each path a type drawn from

\[
  \tau \;::=\; \bot \mid \texttt{int} \mid \texttt{float} \mid \texttt{string}
       \mid \texttt{bool} \mid \mathrm{obj} \mid \mathrm{arr}(\tau)
\]

\noindent
ordered by $\bot \sqsubseteq \tau$ for every $\tau$; $\texttt{int}
\sqsubseteq \texttt{float}$; $\mathrm{arr}(\tau) \sqsubseteq
\mathrm{arr}(\tau')$ whenever $\tau \sqsubseteq \tau'$; and two object types
comparable when their members are, absent members becoming partial. All other
pairs are incomparable.

Nullability is not a point in this lattice. It is a separate flag on the path,
raised by an explicit \texttt{null} or by the member being absent (S4), and by
a null in element position (S16b).

\subsection*{A3. Conditions for acceptance}

An input is accepted, under a configuration, exactly when all of the following
hold.

\begin{enumerate}[label=\textbf{C\arabic*.}]
  \item \textbf{Shape.} The input is an object, or an array all of whose
        elements are objects. \hfill (S0)
  \item \textbf{Unique members.} No object, at any depth, carries the same
        member name twice. \hfill (S10)
  \item \textbf{No directly nested arrays.} No array has an element that is
        itself an array. An array inside an object inside an array is
        permitted. \hfill (S17)
  \item \textbf{Type compatibility.} The types observed at every path have a
        least upper bound in the lattice of A2. \hfill (S16, S18)
  \item \textbf{Map markings are well placed.} Every path named in
        \texttt{maps} exists in the input and is object-valued.
        \hfill (S9)
  \item \textbf{Recursion markings are valid.} For every entry in
        \texttt{recursive}, both paths exist, the second is an ancestor of the
        first, and after inference their shapes are compatible.
        \hfill (S19)
\end{enumerate}

\noindent
C4 is the only condition that requires inference to have run; the rest are
checkable on the input directly.

\subsection*{A4. What is accepted without complaint}

Several inputs that might be expected to fail do not. Each is resolved by
widening or by a collapse rather than by rejection, and the list is worth
stating explicitly because it marks the boundary of C4.

\begin{itemize}
  \item \texttt{\{"a":1\}} and \texttt{\{"a":2.5\}} --- integers and floats
        join at \texttt{float}. \hfill (S1a)
  \item An integer literal too large for OCaml's native \texttt{int} --- the
        column widens to \texttt{float}, with the precision loss recorded in
        S1b rather than raised as an error. \hfill (S1b)
  \item A member present in some documents and absent or \texttt{null} in
        others --- the path becomes partial. \hfill (S4)
  \item A member that is \texttt{null} in every document ---
        \texttt{unit option}. \hfill (S6)
  \item An object that is empty in every document --- a table of bare
        identifiers. \hfill (S8)
  \item An array that is empty in every document --- a table with no value
        column. \hfill (S15)
  \item Objects at one path with different members --- one table whose columns
        are partial. \hfill (S4)
  \item \texttt{[1, null, 2]} --- a nullable element column, not a type
        conflict. \hfill (S16b)
  \item A member name that is not a legal OCaml identifier --- mangled, never
        rejected. \hfill (S11)
\end{itemize}

\subsection*{A5. Rejections}

\begin{center}
\begin{tabular}{@{}llp{6.1cm}@{}}
\hline
\textbf{Condition} & \textbf{Ref} & \textbf{Reported as} \\
\hline
C1 & S0  & The top-level value is not an object or array, or an array element
            is not an object. Names the offending element. \\
C2 & S10 & A repeated member name. Names the path and the key. Per document,
            so one malformed document rejects the input. \\
C3 & S17 & An array element that is an array. Names the path. \\
C4 & S16 & Element types at one array path with no upper bound. Names the path
            and the types. \\
C4 & S18 & Value types at one member path with no upper bound. Names the path,
            the types, and the pair of documents that contributed them. \\
C5 & S9  & A \texttt{maps} entry naming a path that does not exist or is not
            object-valued. \\
C6 & S19 & A \texttt{recursive} entry whose paths do not exist, are not in an
            ancestor relation, or whose shapes are incompatible after
            inference. Names both paths and the mismatch. \\
\hline
\end{tabular}
\end{center}

\subsection*{A6. Status of the condition}

\paragraph{It is decidable, and it is the validator.}
Every condition is decidable on a finite input. C1--C3, C5 and part of C6 are
checkable in a single traversal; C4 and the remainder of C6 require the
per-path type observations that inference computes anyway. The check is
therefore not an additional pass but a reading of inference's own
intermediate state, which is what makes the same condition serviceable as the
ingest-time validator for documents arriving after the schema is fixed.

\paragraph{It is deterministic and structural.}
Given an input and a configuration, acceptance and the resulting schema are
fully determined. No condition depends on a threshold, a sample statistic, or a
heuristic; the two places where such a rule would have been natural --- map
detection (S9) and recursion detection (S19) --- were made configuration
instead, precisely so this holds.

\paragraph{It is not monotone.}
Adding documents to an accepted corpus can make it unacceptable. Everywhere
else in this document, further observations only loosen an inferred type ---
a column widens from \texttt{int} to \texttt{float}, or from total to partial,
and a schema inferred from less data stays valid. C4 is the exception: a
single new document can introduce a type conflict at a path that was uniform
before, and C2 can be violated by one malformed document among many. This is
the practical cost of S16 and S18, and it is recorded here rather than left to
be discovered.

\subsection*{A7. Output}

For an accepted input, Phase~1 produces

\begin{itemize}
  \item a relational schema, with foreign-key constraints materialized and
        nullability recorded per column;
  \item the shredded data;
  \item the generated OCaml modules, one per table, following the two
        principles of S21 --- the record is the row, and the store is
        explicit;
  \item the observation report of S20, giving for every path the counts of
        present-with-value, explicitly-null and absent, which is where the
        distinction S4 collapses out of the type remains visible.
\end{itemize}
